% --- Archon physics formalization source begin ---
% archon:physics
% archon:covers PhyXMiniProblems/problem_phyx_mini_0448.lean
% archon:source-report reports/phyx_mini/problem_phyx_mini_0448.source.json

\chapter{Physics problem phyx\_mini\_0448}
\label{ch:PhyXMiniProblems_problem_phyx_mini_0448}

\paragraph{Problem source.}
\#\# Physical scenario

A \$35 \textbackslash{}, \textbackslash{}text\{ft\}\textasciicircum{}3\$ rigid tank has air at \$225 \textbackslash{}, \textbackslash{}text\{psia\}\$ and ambient \$600 \textbackslash{}, R\$ connected by a valve to a piston/cylinder. The piston of area \$1 \textbackslash{}, \textbackslash{}text\{ft\}\textasciicircum{}2\$ requires \$40 \textbackslash{}, \textbackslash{}text\{psia\}\$ below it to float. The valve is opened, the piston moves slowly \$7 \textbackslash{}, \textbackslash{}text\{ft\}\$ up, and the valve is closed. During the process, air temperature remains at \$600 \textbackslash{}, R\$.

\#\# Question

What is the final pressure in the tank?

\#\# Answer choices

- A: A: 200 psia

- B: B: 212 psia

- C: C: 225 psia

- D: D: 217 psia

\#\# Figure caption (auxiliary; use the image as primary evidence)

The image presents a schematic diagram illustrating a fluid system. 

- **Components and Labels**:
  - **Tank A**: Positioned to the left, represented as a square with the label "Tank A."
  - **Cylinder B**: Shown on the right, depicted as a tall rectangular container with a label "Cylinder B."
  - **Piston**: Inside "Cylinder B," a piston is depicted, signified by a horizontal line with an arrow pointing downwards labeled "Piston."
  - **Valve**: Located at the bottom, between Tank A and Cylinder B, represented by a circle with an X inside, labeled "Valve."
  - **Gravity (g)**: Indicated by an arrow labeled "g" pointing downwards, suggesting gravitational force acting on the piston.

- **Connections and Relationships**:
  - A line connects Tank A and Cylinder B through the valve, suggesting a fluid path.
  - The piston within Cylinder B is influenced by gravity, implying a setup where the piston can move vertically within the cylinder under gravitational force.

- **Additional Details**:
  - There is shading inside the cylinder below the piston, possibly indicating fluid present below the piston.

The diagram likely illustrates a basic hydraulic or pneumatic system where the movement of the piston in Cylinder B is controlled through an open or closed valve and influenced by gravity.

\#\# Dataset metadata

Ideal Gases and Kinetic Theory; Multi-Formula Reasoning, Physical Model Grounding Reasoning

\paragraph{Recorded answer/context.}
D: D: 217 psia

\paragraph{Figure/image path.}
/root/proposal\_for\_physic/science-mango/phyx\_mini\_run/phyx\_data/test\_image/448.png

\paragraph{Formalization target.}
create a compiling Lean file with sorry bodies at `PhyXMiniProblems/problem\_phyx\_mini\_0448.lean`. The Lean declarations must preserve the physical quantities, dimensions or dimensional roles, figure labels, governing-law hypotheses, and final relation expressed by this problem.
Use Mathlib/Physlib names found through LeanExplore where available. If a physics API is missing, introduce faithful local abstractions rather than scalar placeholder aliases.

\begin{theorem}[Physics formalization target]
\label{thm:physics:phyx_mini_0448:target}
\lean{PhyXMiniProblems.ProblemPhyXMini0448.finalTankPressure_eq_217_psia}
\uses{def:physics:phyx-mini-0448:phyxminiproblems-problemphyxmini0448-pressureinpsia, def:physics:phyx-mini-0448:phyxminiproblems-problemphyxmini0448-rigidtankpistonsetup, def:physics:phyx-mini-0448:phyxminiproblems-problemphyxmini0448-matchesproblemandprimaryfigure, def:physics:phyx-mini-0448:phyxminiproblems-problemphyxmini0448-hasphysicalparameters, def:physics:phyx-mini-0448:phyxminiproblems-problemphyxmini0448-satisfiespistonsweepgeometry, def:physics:phyx-mini-0448:phyxminiproblems-problemphyxmini0448-satisfiesisothermalidealgasinventory, def:physics:phyx-mini-0448:phyxminiproblems-problemphyxmini0448-conservesairduringtransfer, def:physics:phyx-mini-0448:phyxminiproblems-problemphyxmini0448-answerpressureinpsia}
The assigned autoformalize agent should translate this physics problem into Lean declarations in the covered file, with theorem and lemma proof bodies written as `by sorry`.
\end{theorem}
\begin{proof}
This is an autoformalization task, not a proof task. Produce statements that can later be proved by the physics prover without weakening the physical model.
\end{proof}

% --- Archon declaration topology begin ---
\section{Declaration topology}
\begin{definition}[Length Quantity]
  \label{def:physics:phyx-mini-0448:phyxminiproblems-problemphyxmini0448-lengthquantity}
  \lean{PhyXMiniProblems.ProblemPhyXMini0448.LengthQuantity}
  A nonnegative physical piston displacement, carrying the length dimension
\end{definition}
\begin{proof}
  This is the declared model definition of Length Quantity.
\end{proof}
\begin{definition}[Volume Quantity]
  \label{def:physics:phyx-mini-0448:phyxminiproblems-problemphyxmini0448-volumequantity}
  \lean{PhyXMiniProblems.ProblemPhyXMini0448.VolumeQuantity}
  A nonnegative physical volume, carrying the dimension L\textasciicircum{}3
\end{definition}
\begin{proof}
  This is the declared model definition of Volume Quantity.
\end{proof}
\begin{definition}[Pressure Quantity]
  \label{def:physics:phyx-mini-0448:phyxminiproblems-problemphyxmini0448-pressurequantity}
  \lean{PhyXMiniProblems.ProblemPhyXMini0448.PressureQuantity}
  A physical absolute pressure
\end{definition}
\begin{proof}
  This is the declared model definition of Pressure Quantity.
\end{proof}
\begin{definition}[Area Quantity]
  \label{def:physics:phyx-mini-0448:phyxminiproblems-problemphyxmini0448-areaquantity}
  \lean{PhyXMiniProblems.ProblemPhyXMini0448.AreaQuantity}
  A nonnegative physical cross-sectional area
\end{definition}
\begin{proof}
  This is the declared model definition of Area Quantity.
\end{proof}
\begin{definition}[pressure In Pascals]
  \label{def:physics:phyx-mini-0448:phyxminiproblems-problemphyxmini0448-pressureinpascals}
  \lean{PhyXMiniProblems.ProblemPhyXMini0448.pressureInPascals}
  \uses{def:physics:phyx-mini-0448:phyxminiproblems-problemphyxmini0448-pressurequantity}
  Read a physical pressure in coherent SI pascals
\end{definition}
\begin{proof}
  This is the declared model definition of pressure In Pascals.
\end{proof}
\begin{definition}[pressure In Psia]
  \label{def:physics:phyx-mini-0448:phyxminiproblems-problemphyxmini0448-pressureinpsia}
  \lean{PhyXMiniProblems.ProblemPhyXMini0448.pressureInPsia}
  \uses{def:physics:phyx-mini-0448:phyxminiproblems-problemphyxmini0448-pressurequantity, def:physics:phyx-mini-0448:phyxminiproblems-problemphyxmini0448-pressureinpascals}
  Read an absolute pressure in pounds per square inch. The denominator is Physlib's dimensionful quantity representing one psi; calling this psia records that all pressures in this problem are absolute, not gauge pressures
\end{definition}
\begin{proof}
  This is the declared model definition of pressure In Psia.
\end{proof}
\begin{definition}[length In Feet]
  \label{def:physics:phyx-mini-0448:phyxminiproblems-problemphyxmini0448-lengthinfeet}
  \lean{PhyXMiniProblems.ProblemPhyXMini0448.lengthInFeet}
  \uses{def:physics:phyx-mini-0448:phyxminiproblems-problemphyxmini0448-lengthquantity}
  Read a piston displacement in feet
\end{definition}
\begin{proof}
  This is the declared model definition of length In Feet.
\end{proof}
\begin{definition}[area In Square Feet]
  \label{def:physics:phyx-mini-0448:phyxminiproblems-problemphyxmini0448-areainsquarefeet}
  \lean{PhyXMiniProblems.ProblemPhyXMini0448.areaInSquareFeet}
  \uses{def:physics:phyx-mini-0448:phyxminiproblems-problemphyxmini0448-areaquantity}
  Read a cross-sectional area in square feet
\end{definition}
\begin{proof}
  This is the declared model definition of area In Square Feet.
\end{proof}
\begin{definition}[volume In Cubic Feet]
  \label{def:physics:phyx-mini-0448:phyxminiproblems-problemphyxmini0448-volumeincubicfeet}
  \lean{PhyXMiniProblems.ProblemPhyXMini0448.volumeInCubicFeet}
  \uses{def:physics:phyx-mini-0448:phyxminiproblems-problemphyxmini0448-volumequantity}
  Read a physical volume using feet as the selected length unit
\end{definition}
\begin{proof}
  This is the declared model definition of volume In Cubic Feet.
\end{proof}
\begin{definition}[temperature In Rankine]
  \label{def:physics:phyx-mini-0448:phyxminiproblems-problemphyxmini0448-temperatureinrankine}
  \lean{PhyXMiniProblems.ProblemPhyXMini0448.temperatureInRankine}
  Rankine is the absolute Fahrenheit scale. Physlib's Temperature deliberately uses an arbitrary absolute unit; in this setup that unit is fixed to TemperatureUnit.absoluteFahrenheit, and this is its real-valued readout
\end{definition}
\begin{proof}
  This is the declared model definition of temperature In Rankine.
\end{proof}
\begin{definition}[Gas Species]
  \label{def:physics:phyx-mini-0448:phyxminiproblems-problemphyxmini0448-gasspecies}
  \lean{PhyXMiniProblems.ProblemPhyXMini0448.GasSpecies}
  The gas species named by the problem
\end{definition}
\begin{proof}
  This is the declared model definition of Gas Species.
\end{proof}
\begin{definition}[Vessel]
  \label{def:physics:phyx-mini-0448:phyxminiproblems-problemphyxmini0448-vessel}
  \lean{PhyXMiniProblems.ProblemPhyXMini0448.Vessel}
  The two vessels labelled in the primary image
\end{definition}
\begin{proof}
  This is the declared model definition of Vessel.
\end{proof}
\begin{definition}[Tank Behavior]
  \label{def:physics:phyx-mini-0448:phyxminiproblems-problemphyxmini0448-tankbehavior}
  \lean{PhyXMiniProblems.ProblemPhyXMini0448.TankBehavior}
  The stated mechanical character of tank A
\end{definition}
\begin{proof}
  This is the declared model definition of Tank Behavior.
\end{proof}
\begin{definition}[Piston Motion Regime]
  \label{def:physics:phyx-mini-0448:phyxminiproblems-problemphyxmini0448-pistonmotionregime}
  \lean{PhyXMiniProblems.ProblemPhyXMini0448.PistonMotionRegime}
  The slow piston motion assumed by the textbook model
\end{definition}
\begin{proof}
  This is the declared model definition of Piston Motion Regime.
\end{proof}
\begin{definition}[Thermal Regime]
  \label{def:physics:phyx-mini-0448:phyxminiproblems-problemphyxmini0448-thermalregime}
  \lean{PhyXMiniProblems.ProblemPhyXMini0448.ThermalRegime}
  The thermal constraint stated for the valve-opening process
\end{definition}
\begin{proof}
  This is the declared model definition of Thermal Regime.
\end{proof}
\begin{definition}[Air Transfer Direction]
  \label{def:physics:phyx-mini-0448:phyxminiproblems-problemphyxmini0448-airtransferdirection}
  \lean{PhyXMiniProblems.ProblemPhyXMini0448.AirTransferDirection}
  Direction in which air crosses the open valve
\end{definition}
\begin{proof}
  This is the declared model definition of Air Transfer Direction.
\end{proof}
\begin{definition}[Valve Event]
  \label{def:physics:phyx-mini-0448:phyxminiproblems-problemphyxmini0448-valveevent}
  \lean{PhyXMiniProblems.ProblemPhyXMini0448.ValveEvent}
  Three events in the valve operation described in the prose
\end{definition}
\begin{proof}
  This is the declared model definition of Valve Event.
\end{proof}
\begin{definition}[Valve Position]
  \label{def:physics:phyx-mini-0448:phyxminiproblems-problemphyxmini0448-valveposition}
  \lean{PhyXMiniProblems.ProblemPhyXMini0448.ValvePosition}
  Whether the valve is open or closed at a given event
\end{definition}
\begin{proof}
  This is the declared model definition of Valve Position.
\end{proof}
\begin{definition}[Vertical Direction]
  \label{def:physics:phyx-mini-0448:phyxminiproblems-problemphyxmini0448-verticaldirection}
  \lean{PhyXMiniProblems.ProblemPhyXMini0448.VerticalDirection}
  Vertical direction of the piston motion or gravity arrow
\end{definition}
\begin{proof}
  This is the declared model definition of Vertical Direction.
\end{proof}
\begin{definition}[Figure Component]
  \label{def:physics:phyx-mini-0448:phyxminiproblems-problemphyxmini0448-figurecomponent}
  \lean{PhyXMiniProblems.ProblemPhyXMini0448.FigureComponent}
  Physical components distinguished in image 448.png
\end{definition}
\begin{proof}
  This is the declared model definition of Figure Component.
\end{proof}
\begin{definition}[Figure Label]
  \label{def:physics:phyx-mini-0448:phyxminiproblems-problemphyxmini0448-figurelabel}
  \lean{PhyXMiniProblems.ProblemPhyXMini0448.FigureLabel}
  Text or symbols printed in the primary image
\end{definition}
\begin{proof}
  This is the declared model definition of Figure Label.
\end{proof}
\begin{definition}[Horizontal Placement]
  \label{def:physics:phyx-mini-0448:phyxminiproblems-problemphyxmini0448-horizontalplacement}
  \lean{PhyXMiniProblems.ProblemPhyXMini0448.HorizontalPlacement}
  Relative left-to-right placement visible in the schematic
\end{definition}
\begin{proof}
  This is the declared model definition of Horizontal Placement.
\end{proof}
\begin{definition}[Tank Piston Figure]
  \label{def:physics:phyx-mini-0448:phyxminiproblems-problemphyxmini0448-tankpistonfigure}
  \lean{PhyXMiniProblems.ProblemPhyXMini0448.TankPistonFigure}
  \uses{def:physics:phyx-mini-0448:phyxminiproblems-problemphyxmini0448-vessel, def:physics:phyx-mini-0448:phyxminiproblems-problemphyxmini0448-verticaldirection, def:physics:phyx-mini-0448:phyxminiproblems-problemphyxmini0448-figurecomponent, def:physics:phyx-mini-0448:phyxminiproblems-problemphyxmini0448-figurelabel, def:physics:phyx-mini-0448:phyxminiproblems-problemphyxmini0448-horizontalplacement}
  Structured transcription of the components and geometry in image 448.png
\end{definition}
\begin{proof}
  This is the declared model definition of Tank Piston Figure.
\end{proof}
\begin{definition}[Rigid Tank Piston Setup]
  \label{def:physics:phyx-mini-0448:phyxminiproblems-problemphyxmini0448-rigidtankpistonsetup}
  \lean{PhyXMiniProblems.ProblemPhyXMini0448.RigidTankPistonSetup}
  \uses{def:physics:phyx-mini-0448:phyxminiproblems-problemphyxmini0448-lengthquantity, def:physics:phyx-mini-0448:phyxminiproblems-problemphyxmini0448-volumequantity, def:physics:phyx-mini-0448:phyxminiproblems-problemphyxmini0448-pressurequantity, def:physics:phyx-mini-0448:phyxminiproblems-problemphyxmini0448-areaquantity, def:physics:phyx-mini-0448:phyxminiproblems-problemphyxmini0448-gasspecies, def:physics:phyx-mini-0448:phyxminiproblems-problemphyxmini0448-tankbehavior, def:physics:phyx-mini-0448:phyxminiproblems-problemphyxmini0448-pistonmotionregime, def:physics:phyx-mini-0448:phyxminiproblems-problemphyxmini0448-thermalregime, def:physics:phyx-mini-0448:phyxminiproblems-problemphyxmini0448-airtransferdirection, def:physics:phyx-mini-0448:phyxminiproblems-problemphyxmini0448-valveevent, def:physics:phyx-mini-0448:phyxminiproblems-problemphyxmini0448-valveposition, def:physics:phyx-mini-0448:phyxminiproblems-problemphyxmini0448-tankpistonfigure}
  Independent physical quantities and measured scalar components of the setup. In particular, tankAFinalPressure is an unconstrained physical-pressure field: it is not defined from 217 or from an answer choice
\end{definition}
\begin{proof}
  This is the declared model definition of Rigid Tank Piston Setup.
\end{proof}
\begin{definition}[Matches Problem And Primary Figure]
  \label{def:physics:phyx-mini-0448:phyxminiproblems-problemphyxmini0448-matchesproblemandprimaryfigure}
  \lean{PhyXMiniProblems.ProblemPhyXMini0448.MatchesProblemAndPrimaryFigure}
  \uses{def:physics:phyx-mini-0448:phyxminiproblems-problemphyxmini0448-pressureinpsia, def:physics:phyx-mini-0448:phyxminiproblems-problemphyxmini0448-lengthinfeet, def:physics:phyx-mini-0448:phyxminiproblems-problemphyxmini0448-areainsquarefeet, def:physics:phyx-mini-0448:phyxminiproblems-problemphyxmini0448-volumeincubicfeet, def:physics:phyx-mini-0448:phyxminiproblems-problemphyxmini0448-temperatureinrankine, def:physics:phyx-mini-0448:phyxminiproblems-problemphyxmini0448-rigidtankpistonsetup}
  Numerical data from the prose and qualitative data from the primary image. No field here fixes the requested final tank pressure, 217 psia, or an answer-choice label
\end{definition}
\begin{proof}
  This is the declared model definition of Matches Problem And Primary Figure.
\end{proof}
\begin{definition}[Has Physical Parameters]
  \label{def:physics:phyx-mini-0448:phyxminiproblems-problemphyxmini0448-hasphysicalparameters}
  \lean{PhyXMiniProblems.ProblemPhyXMini0448.HasPhysicalParameters}
  \uses{def:physics:phyx-mini-0448:phyxminiproblems-problemphyxmini0448-pressureinpsia, def:physics:phyx-mini-0448:phyxminiproblems-problemphyxmini0448-lengthinfeet, def:physics:phyx-mini-0448:phyxminiproblems-problemphyxmini0448-areainsquarefeet, def:physics:phyx-mini-0448:phyxminiproblems-problemphyxmini0448-volumeincubicfeet, def:physics:phyx-mini-0448:phyxminiproblems-problemphyxmini0448-temperatureinrankine, def:physics:phyx-mini-0448:phyxminiproblems-problemphyxmini0448-rigidtankpistonsetup}
  Positivity conditions selecting the intended physical branch
\end{definition}
\begin{proof}
  This is the declared model definition of Has Physical Parameters.
\end{proof}
\begin{definition}[Satisfies Piston Sweep Geometry]
  \label{def:physics:phyx-mini-0448:phyxminiproblems-problemphyxmini0448-satisfiespistonsweepgeometry}
  \lean{PhyXMiniProblems.ProblemPhyXMini0448.SatisfiesPistonSweepGeometry}
  \uses{def:physics:phyx-mini-0448:phyxminiproblems-problemphyxmini0448-lengthinfeet, def:physics:phyx-mini-0448:phyxminiproblems-problemphyxmini0448-areainsquarefeet, def:physics:phyx-mini-0448:phyxminiproblems-problemphyxmini0448-volumeincubicfeet, def:physics:phyx-mini-0448:phyxminiproblems-problemphyxmini0448-rigidtankpistonsetup}
  The vertical piston displacement sweeps out area * rise beneath the piston. This geometric law contains no pressure value and no requested conclusion
\end{definition}
\begin{proof}
  This is the declared model definition of Satisfies Piston Sweep Geometry.
\end{proof}
\begin{definition}[Satisfies Isothermal Ideal Gas Inventory]
  \label{def:physics:phyx-mini-0448:phyxminiproblems-problemphyxmini0448-satisfiesisothermalidealgasinventory}
  \lean{PhyXMiniProblems.ProblemPhyXMini0448.SatisfiesIsothermalIdealGasInventory}
  \uses{def:physics:phyx-mini-0448:phyxminiproblems-problemphyxmini0448-pressureinpsia, def:physics:phyx-mini-0448:phyxminiproblems-problemphyxmini0448-volumeincubicfeet, def:physics:phyx-mini-0448:phyxminiproblems-problemphyxmini0448-temperatureinrankine, def:physics:phyx-mini-0448:phyxminiproblems-problemphyxmini0448-rigidtankpistonsetup}
  Ideal-gas inventory equations in one consistent mixed-unit system. The first two equations describe the air remaining in the same rigid tank before and after transfer. The third describes the admitted air occupying the volume swept beneath a quasistatic piston held at its floating pressure. A single temperature field expresses the stated isothermal condition
\end{definition}
\begin{proof}
  This is the declared model definition of Satisfies Isothermal Ideal Gas Inventory.
\end{proof}
\begin{definition}[Conserves Air During Transfer]
  \label{def:physics:phyx-mini-0448:phyxminiproblems-problemphyxmini0448-conservesairduringtransfer}
  \lean{PhyXMiniProblems.ProblemPhyXMini0448.ConservesAirDuringTransfer}
  \uses{def:physics:phyx-mini-0448:phyxminiproblems-problemphyxmini0448-rigidtankpistonsetup}
  Air inventory is conserved while the valve is open: the loss from rigid tank A equals the amount admitted to cylinder B. This law contains no pressure answer
\end{definition}
\begin{proof}
  This is the declared model definition of Conserves Air During Transfer.
\end{proof}
\begin{definition}[Answer Choice]
  \label{def:physics:phyx-mini-0448:phyxminiproblems-problemphyxmini0448-answerchoice}
  \lean{PhyXMiniProblems.ProblemPhyXMini0448.AnswerChoice}
  Labels of the four answers printed with the problem
\end{definition}
\begin{proof}
  This is the declared model definition of Answer Choice.
\end{proof}
\begin{definition}[answer Pressure In Psia]
  \label{def:physics:phyx-mini-0448:phyxminiproblems-problemphyxmini0448-answerpressureinpsia}
  \lean{PhyXMiniProblems.ProblemPhyXMini0448.answerPressureInPsia}
  \uses{def:physics:phyx-mini-0448:phyxminiproblems-problemphyxmini0448-answerchoice}
  Absolute-pressure value printed beside each displayed answer, in psia
\end{definition}
\begin{proof}
  This is the declared model definition of answer Pressure In Psia.
\end{proof}
\begin{definition}[recorded Answer Choice]
  \label{def:physics:phyx-mini-0448:phyxminiproblems-problemphyxmini0448-recordedanswerchoice}
  \lean{PhyXMiniProblems.ProblemPhyXMini0448.recordedAnswerChoice}
  \uses{def:physics:phyx-mini-0448:phyxminiproblems-problemphyxmini0448-answerchoice}
  The answer label recorded in the supplied dataset
\end{definition}
\begin{proof}
  This is the declared model definition of recorded Answer Choice.
\end{proof}
% --- Archon declaration topology end ---
% --- Archon physics formalization source end ---
